% Volume IX: Institutional Design
% Part XVII. Designing for Delayed Truth
% Chapter 97: Reopenable Records
% Spine: proof-spines/ch097-spine.md

\chapter{Reopenable Records}
\label{ch:ch097}

Institutions should preserve not only what was decided, but the provenance, uncertainty, dissent, and dependency structure that made the decision possible. Under conditions of delayed truth, a record that stores only its terminal conclusion invites later reviewers to mistake administrative closure for epistemic completion. A reopenable record therefore keeps track of excluded alternatives as well as accepted ones, and it states why those alternatives were set aside at the time.

The point is not to abolish decision, but to ensure that later evidence has somewhere responsible to attach. If provenance has been erased, if uncertainty has been flattened, if dissent has been dropped, or if dependencies among premises have been hidden, correction becomes needlessly difficult. Reopenable records are thus archives designed for revision: they preserve the materials by which institutional judgment can be updated without pretending that the first account was final merely because it was first.
